\scanpage{141}\mspage{15}
% archive sheet number: 140
dans (2.5.) n'est autre que l'inclusion canonique d'une somme amalgamée partielle, indexée par $\Delta(C_0,f)$, dans la somme am. \unsure{correspondant|correspondante} à l'ens. d'indices plus grand. La fonctorialité de cette flèche $u_1(f)$, pour $f$ variable, ``est immédiate'', elle résulte de la commutativité du carré, associé à un \deleted{flèche} morphisme \inserted{$\unsure{v|\nu|\alpha}=(v_X,v_Y)$} de flèches
\diagram{\begin{tikzcd}[ampersand replacement=\&,column sep=large,row sep=large,execute at end picture={\node[anchor=east] at ([xshift=-35pt]current bounding box.west) {(2.\unsure{6|8}.)};}]
X\arrow[r,dash,"f"]\arrow[d,dash,"\deleted{\unsure{u|\alpha}}v_X"']\&Y\arrow[d,dash,"v_Y"]\\
X'\arrow[r,dash,"f'"']\&Y'
\end{tikzcd}}
, \deleted{\unsure{dans|deux}} \inserted{sur le} diagramme d'applications entre ens. d'indices
\diagram{\begin{tikzcd}[ampersand replacement=\&,column sep=huge,row sep=huge,execute at end picture={\node[anchor=east] at ([xshift=-35pt]current bounding box.west) {(2.7)};}]
\Delta(C_0,f)\arrow[r,hook,"u_1(f)"]\arrow[d,"\substack{\text{\unsure{dép.|dep.} \unsure{fonctlle|fonct.}}\\\text{en }f\text{ de }\Delta(C_0,f)}"']\&\Delta(C_1,f)\arrow[d,"\substack{\text{dép. fonct. en }f\\\text{de }\Delta(C_1,f)}"]\\
\Delta(C_0,f')\arrow[r,hook,"u_1(f')"]\&\Delta(C_1,f')
\end{tikzcd}}
(\deleted{\unsure{ce|la}} \inserted{commutativité} \deleted{\unsure{ce}} qui \unsure{revient|revîent} à dire que $\Delta(C_0,f)$ peut être considéré comme un foncteur en le couple $\unsure{C_0,f|C_0\ f}$ \unsure{\&|et} comme objet de la catégorie produit $\uline{\unsure{\mathscr P|\mathcal P}_{\mathrm{pet.}}(\mathrm{Fl}_{\mathrm{coquar}}(M))}\times\uline{\mathrm{Fl}}(M)$\ldots). Le carré commutatif (2.7) est la partie ``\unsure{indices|indice}'' d'un carré commutatif entre familles d'objets $\unsure{\mathscr X|\mathcal X|\mathfrak X}(C_0,f)$, $\mathscr X(C_0,f')$, $\mathscr X(C_1,f)$, $\mathscr X(C_1,f')$
